\section{Introduction}

Temporal representation quality under drift is not exhausted by endpoint task
accuracy. A model can retain acceptable pre-drift performance and still require
substantially different post-drift re-learning effort from another model with
similar static accuracy. The empirical question is therefore not only whether a
representation is accurate, but how costly it is to re-tune once the signal
changes.

The operational object in this paper is the minimum realizable adaptation burden
over protocol-window choices,
\[
  R^*(m,\pi) \coloneqq \min_w R(m,\pi,w).
\]
Here $m$ denotes the trained model family and $\pi$ denotes the adaptation
protocol. This object connects the empirical program to the temporal-validity
functional $V(\phi,\tau)$: inferential utility curves over horizon identify how
useful present representations remain for future targets, while $R^*$ measures
the downstream cost of restoring utility after drift.

The central empirical claim is hierarchical rather than universal. Effective
spectral rank separates degradation regimes across architectural families.
Within a recurrent local regime, activation velocity $\Delta_v$ predicts burden
more strongly than effective rank, even though effective rank still varies over
a wide range. Those two facts imply that family choice and within-family scale
choice move burden through distinct empirical channels.

That hierarchy matters because the literature often compresses stability,
plasticity, representational geometry, and memory into a single tradeoff. The
results below instead identify two observational levels. Family-level geometry
determines which burden regime a model enters. Within at least one recurrent
family, dynamic reactivity determines where the model sits inside that regime.
